\uuid{QGzL}
\titre{Trigonométrie}
\auteur{bodin ; hanani ; mzari}
\organisation{exo7}
\licence{CC BY-NC-SA 4.0}
\datecreate{2020-01-20}
\niveau{L1}
\module{Analyse}
\chapitre{Fonctions circulaires et hyperboliques inverses}
\sousChapitre{Autre}
\difficulte{3}
\theme{Fonctions usuelles, trigonométrie}
\exo7qcmid{lille-1.fonctions-usuelles.026}

\contenu{
\begin{question}
\qclassification{126.00}
\qtags{motcle=trigonométrie}
Soit $f(x)= \arccos (\frac{1-x^2}{1+x^2})$. On notera $D_f$ le domaine de définition de $f$. Quelles sont les assertions vraies ?
\begin{answers}
    \good{$D_f= \Rr$}

    \bad{$D_f=[-1,1]$}

    \good{$ f(x)= 2\arctan x, \, \forall x\ge 0$}

    \good{$f(x)  = 2\arctan |x|, \, \forall x \in \Rr$}
\end{answers}
\begin{explanations}
$f$ est définie si $-1\le \frac{1-x^2}{1+x^2} \le 1 $, ce qui est le cas pour tout $x \in \Rr$.
$f$ est dérivable sur $\Rr^*$ et $f'(x) = \frac{2x}{|x|(1+x^2)} =\left\{\begin{array}{cc}\frac{2}{1+x^2},& \mbox{si} \, x>0 \\ -\frac{2}{1+x^2},& \mbox{si} \, x<0 \end{array}\right. $.
On déduit que
$f(x)= \left\{\begin{array}{cc}2\arctan x + c_1, & \mbox{si} \, x\ge 0 \\ -2\arctan x + c_2, & \mbox{si} \, x\le 0 \end{array}\right. $, où $c_1$ et $ c_2$  sont des constantes.
Comme $f$ est continue en $0$ (ou par identification en $+\infty$ et en $-\infty$), on obtient $c_1=c_2=0$. Ainsi
$f(x)= \left\{\begin{array}{cc}2\arctan x, & \mbox{si} \, x\ge 0 \\ -2\arctan x, & \mbox{si} \, x\le 0 \end{array}\right.$.  Donc $f(x)  = 2\arctan |x|, \, \forall x \in \Rr$.

\end{explanations}
\end{question}
}
